\section*{Chapitre \ref{ch:g11:circuits-and-power} --- Circuits électriques et puissance}

\begin{solution}{exo:g11:circuits-and-power:1}
$q = It = 2.0 \times 90 \times 60 = \qty{1.1e4}{C}$. Nombre de charges
élémentaires~: $q/e = 1.08\times10^{4}/1.6\times10^{-19} \approx
\num{6.8e22}$.
\end{solution}

\begin{solution}{exo:g11:circuits-and-power:2}
$I = U/R = 12/220 = \qty{0.055}{A} = \qty{55}{mA}$~;
$R = U/I = 230/0.50 = \qty{460}{\ohm}$.
\end{solution}

\begin{solution}{exo:g11:circuits-and-power:3}
Série~: $120 + 60 = \qty{180}{\ohm}$. Parallèle~:
$\frac{120 \times 60}{180} = \qty{40}{\ohm}$ --- plus petite que les
deux~: même \omterm{def:g11:circuits-and-power:voltage}{tension}, mais le courant a une
seconde voie, donc plus d'intensité totale et moins de
\omterm{def:g11:circuits-and-power:resistance}{résistance}.
\end{solution}

\begin{solution}{exo:g11:circuits-and-power:4}
$P = UI = 230 \times 8.7 = \qty{2.0}{kW}$. En $t = \qty{240}{s}$~:
$E = Pt = 2000 \times 240 = \qty{4.8e5}{J} = 4.8\times10^{5}/3.6
\times10^{6} \approx \qty{0.13}{kW.h}$.
\end{solution}

\begin{solution}{exo:g11:circuits-and-power:5}
$E = 60 \times 5.0 = \qty{300}{W.h} = \qty{0.30}{kW.h} = 0.30 \times
3.6\times10^{6} = \qty{1.1e6}{J}$. Coût~: $0.30 \times 0.25 = 0.075$
euro par jour.
\end{solution}

\begin{solution}{exo:g11:circuits-and-power:6}
\emph{1.} $R = U^2/P = 230^2/1200 = \qty{44}{\ohm}$~;
$I = P/U = 1200/230 = \qty{5.2}{A}$.

\emph{2.} $P' = U'^2/R = 115^2/44.1 = \qty{300}{W}$~: un quart, pas la
moitié, car diviser $U$ par deux divise aussi $I$ par deux, et
$P = UI \propto U^2$.
\end{solution}

\begin{solution}{exo:g11:circuits-and-power:7}
$R_{\text{s}} = \qty{60}{\ohm}$, donc $I = 12/60 = \qty{0.20}{A}$.
$U_1 = R_1 I = \qty{8.0}{V}$, $U_2 = \qty{4.0}{V}$ (somme \qty{12}{V})~;
$P_1 = U_1 I = \qty{1.6}{W}$, $P_2 = \qty{0.8}{W}$~: total \qty{2.4}{W}
$= UI = 12 \times 0.20$, comme il se doit.
\end{solution}

\begin{solution}{exo:g11:circuits-and-power:8}
$I_1 = 12/60 = \qty{0.20}{A}$, $I_2 = 12/30 = \qty{0.40}{A}$, total
$I = \qty{0.60}{A}$. Équivalente~: $R = U/I = 12/0.60 = \qty{20}{\ohm}$,
et en effet $\frac{60 \times 30}{60 + 30} = \qty{20}{\ohm}$.
\end{solution}

\begin{solution}{exo:g11:circuits-and-power:9}
$\mathcal{E} - 0.50\,r = 5.5$ et $\mathcal{E} - 2.0\,r = 4.0$~;
en soustrayant, $1.5\,r = 1.5$, donc $r = \qty{1.0}{\ohm}$ et
$\mathcal{E} = \qty{6.0}{V}$. Court-circuit~: $\mathcal{E}/r =
\qty{6.0}{A}$.
\end{solution}

\begin{solution}{exo:g11:circuits-and-power:10}
Bouilloire $2200/230 = \qty{9.6}{A}$, grille-pain $1000/230 =
\qty{4.3}{A}$~: total \qty{13.9}{A} $< \qty{16}{A}$, le
\omterm{def:g11:circuits-and-power:joule}{fusible} tient. Radiateur $1500/230 =
\qty{6.5}{A}$~: total \qty{20.4}{A} $> \qty{16}{A}$, il fond
(équivalemment, $\qty{4.7}{kW} > 230 \times 16 \approx \qty{3.7}{kW}$).
\end{solution}

\begin{solution}{exo:g11:circuits-and-power:11}
À \qty{200}{V}~: $I = 10^4/200 = \qty{50}{A}$, perte $RI^2 = 2.0 \times
50^2 = \qty{5.0}{kW}$ --- la moitié de la livraison. À \qty{4.0}{kV}~:
$I = \qty{2.5}{A}$, perte $2.0 \times 2.5^2 = \qty{12.5}{W}$. Rapport
$400 = 20^2$~: vingt fois la
\omterm{def:g11:circuits-and-power:voltage}{tension} signifie vingt fois moins
d'intensité, et la perte va en $I^2$.
\end{solution}

\begin{solution}{exo:g11:circuits-and-power:12}
Quatre montages~: tous en série, $3 \times 60 = \qty{180}{\ohm}$~; tous
en parallèle, $60/3 = \qty{20}{\ohm}$~; un en série avec deux en
parallèle, $60 + 30 = \qty{90}{\ohm}$~; un en parallèle avec deux en
série, $\frac{60 \times 120}{180} = \qty{40}{\ohm}$.
\end{solution}

\begin{solution}{exo:g11:circuits-and-power:13}
\emph{1.} $I = \mathcal{E}/(R + r) = 4.5/7.5 = \qty{0.60}{A}$~;
$U = RI = \qty{3.6}{V}$~; utile $P = UI = \qty{2.2}{W}$~; perdue $rI^2 =
1.5 \times 0.60^2 = \qty{0.54}{W}$~; $\eta = U/\mathcal{E} = 3.6/4.5 =
0.80$.

\emph{2.} Par l'indication, $P = \mathcal{E}^2 \big/ \big[(R - r)^2/R +
4r\big]$~; le crochet est minimal ($= 4r$) exactement lorsque $R = r$,
donc $P_{\max} = \mathcal{E}^2/(4r) = 4.5^2/6.0 = \qty{3.4}{W}$.
\end{solution}

\begin{solution}{exo:g11:circuits-and-power:14}
Chaleur~: $Q = 1.0 \times 4180 \times (100 - 20) = \qty{3.3e5}{J}$.
Électrique~: $E = Q/0.90 = \qty{3.7e5}{J}$. Temps~: $t = E/P =
3.7\times10^{5}/2200 \approx \qty{170}{s}$, moins de trois minutes.
Coût~: $3.7\times10^{5}/3.6\times10^{6} \approx \qty{0.10}{kW.h}$,
environ $2{,}6$~centimes.
\end{solution}

\begin{solution}{exo:g11:circuits-and-power:15}
\emph{1.} $q = 60 \times 3600 = \qty{2.2e5}{C}$~; $E \approx \mathcal{E}q
= 12.7 \times 2.16\times10^{5} = \qty{2.7e6}{J} \approx \qty{0.76}{kW.h}$.

\emph{2.} $U = 12.7 - 0.020 \times 150 = \qty{9.7}{V}$~; délivrée
$UI = \qty{1.5}{kW}$~; perdue $rI^2 = 0.020 \times 150^2 = \qty{450}{W}$.

\emph{3.} Les phares sont aux bornes~: pendant le démarrage, la
\omterm{def:g11:circuits-and-power:voltage}{tension} aux bornes chute de $12.7$ à
$\qty{9.7}{V}$, donc ils reçoivent moins de
\omterm{def:g10:energy-conservation:power}{puissance} et baissent.
\end{solution}

\begin{solution}{pb:g11:circuits-and-power:1}
\textbf{1.} Chaque seconde, $E_p = mgh = 5.0\times10^{4} \times 9.81
\times 200 = \qty{9.8e7}{J}$~: \omterm{def:g10:energy-conservation:power}{puissance}
disponible \qty{98}{MW}.

\textbf{2.} $0.90 \times 98 = \qty{88}{MW}$.

\textbf{3.} $I = P/U = 8.8\times10^{7}/2.0\times10^{4} = \qty{4.4e3}{A}$.

\textbf{4.} $R_\ell I^2 = 32 \times (4.4\times10^{3})^2 \approx
\qty{6.2e8}{W} = \qty{620}{MW}$ --- sept fois la sortie de la centrale~:
transmettre à \qty{20}{kV} est impossible.

\textbf{5.} Avec $P = UI$ fixe, seule l'élévation de $U$ diminue $I$~:
un transformateur.

\textbf{6.} $I = 8.8\times10^{7}/4.0\times10^{5} = \qty{220}{A}$.

\textbf{7.} $R_\ell I^2 = 32 \times 220^2 \approx \qty{1.6}{MW}$.

\textbf{8.} $1.6/88 \approx 0.018$~: environ $1{,}8\,\%$.

\textbf{9.} Fraction $= R_\ell I^2/P = R_\ell (P/U)^2/P = R_\ell P/U^2
\propto 1/U^2$~: de \qty{20}{kV} à \qty{400}{kV} elle chute d'un
facteur $20^2 = 400$, et en effet $\qty{620}{MW}/400 \approx
\qty{1.6}{MW}$.

\textbf{10.} $R_\ell I = 32 \times 220 \approx \qty{7.1}{kV}$~: $1{,}8\,\%$
de \qty{400}{kV}, la même fraction que la question~8, puisque
$R_\ell I/U = R_\ell I^2/UI$.

\textbf{11.} $\qty{1}{kW.h} = \qty{3.6e6}{J}$~: le compteur multiplie la
\omterm{def:g10:energy-conservation:power}{puissance} tirée par le temps pendant lequel
elle est tirée, et additionne.

\textbf{12.} $I = 900/230 = \qty{3.9}{A}$~;
$R = U^2/P = 230^2/900 \approx \qty{59}{\ohm}$.

\textbf{13.} $E = 900 \times 180 = \qty{1.6e5}{J} = \qty{0.045}{kW.h}$~:
environ $1{,}1$~centime.

\textbf{14.} La physique de $RI^2$ est identique~; seul le but
diffère --- sur la ligne la chaleur est indésirable, dans le
grille-pain la chaleur \emph{est} le produit.

\textbf{15.} $3.9 + 9.6 = \qty{13.5}{A} < \qty{16}{A}$~: tient.
Radiateur $2000/230 = \qty{8.7}{A}$~: total $\qty{22.2}{A} >
\qty{16}{A}$, le \omterm{def:g11:circuits-and-power:joule}{fusible} fond.

\textbf{16.} $0.05 \times 13.5^2 = \qty{9.1}{W}$~: $9.1/3100 \approx
0.3\,\%$ de la \omterm{def:g10:energy-conservation:power}{puissance} livrée --- négligeable
sur des fils courts.

\textbf{17.} $\eta = 0.90 \times 0.99 \times 0.982 \times 0.99 \times
0.997 \approx 0.86$.

\textbf{18.} Un \omterm{def:g10:energy-conservation:kwh}{kilowatt-heure} toasté demande
$3.6\times10^{6}/0.86 \approx \qty{4.2e6}{J}$ en amont~:
$m = E/(gh) = 4.2\times10^{6}/(9.81 \times 200) \approx \qty{2.1e3}{kg}$,
environ \qty{2.1}{m^3} d'eau.

\textbf{19.} $R_\ell P/U^2 = 32 \times 8.8\times10^{7}/(2.0\times
10^{4})^2 \approx 7$. Une «~fraction~» supérieure à~$1$ signifie que la
ligne dissiperait plus qu'elle ne transporte~: la livraison ne peut
pas se faire du tout à cette
\omterm{def:g11:circuits-and-power:voltage}{tension}.

\textbf{20.} Un \omterm{def:g10:energy-conservation:kwh}{kilowatt-heure} toasté coûte à la
montagne environ deux tonnes --- \qty{2.1}{m^3} --- d'eau tombant de
\qty{200}{m}, et quelque $14\,\%$ de
l'\omterm{def:g10:energy-conservation:energy}{énergie} meurt en route, surtout dans les
turbines et transformateurs (moins de $2\,\%$ sur \qty{1000}{km} de
ligne). L'unique tour qui rend le voyage possible est le transformateur~:
élever à \qty{400}{kV} divise l'intensité par $20$ et la perte Joule par
$400$.
\end{solution}
